\documentclass[letter,11pt]{article}
\usepackage{hyperref}
\usepackage{graphicx}
\usepackage{fancyhdr}
\usepackage{multicol}
\pagestyle{fancy}
\usepackage[letterpaper, margin=1in]{geometry}
\geometry{letterpaper}
\usepackage{parskip} % Disable initial indent
\usepackage{attachfile}

\usepackage[utf8]{inputenc}
\fancyhf{}
\renewcommand{\headrulewidth}{0pt} % Remove default underline from header package
\rhead{CYBR/CMSC 691 - Lab 4}
%\rhead{}
\lhead{\begin{picture}(0,0) \put(0,-10){\includegraphics[width=1.1cm]{Images/UMBC-vertical}} \end{picture}}
\cfoot{\thepage}
\rfoot{Spring 2020}
\lfoot{Lab 4}
%\AtEndDocument{\vfill \hfill \LaTeX}

\begin{document}

\section*{Lab 4 - Domain-Specific Features}
\paragraph{}Included in the virtual machine is a collection of malware which is part of Endgame's Ember dataset. Using the provided code, build a model based on the training data, and evaluate it on the testing data. Study how the change in features changes the model's performance. The \texttt{trainModel.py} script trains both a Logistic Regression model and a Decision Tree. They are different algorithms which attempt to do the same thing: develop a model to represent the data observed.

\textbf{Points:} 20

\section*{Instructions}
\begin{enumerate}
    \item Write a script to break the data into separate training and testing directories, with an 80\%/20\% split. You should use symbolic links instead of copying the files, since it will save disk space. In Python: \texttt{os.symlink(existingFilePath, newLinkPath)}.
    \item Check out \texttt{pefileInfo.py}\attachfile{Lab4/pefileInfo.py}. It shows basic usage of the \texttt{pefile} module in Python, and how to get detailed information out from an individual sample. This is a place to get acquanted with the \texttt{pefile} module, and learn how to extract desired information. You'll want to refer to the documentation at \url{https://github.com/erocarrera/pefile/blob/wiki/UsageExamples.md} to get started.
    
    \item Using the \texttt{makeDataset.py} \attachfile{Lab4/makeDataset.py} script, copy \& paste some of the features you've found to the \texttt{featurizer()} function, at the end where the function returns a list.
    
    \item Using the \texttt{trainModel.py} \attachfile{Lab4/trainModel.py} script, train the models with the dataset you've created.
    
    \item Make changes to the features, re-create the dataset, and retrain the models. Observe how features improve and worsen the models' accuracy.
    \item Don't forget to include screenshots.
\end{enumerate}

\section*{Questions}
\paragraph{}In a separate document, provide your answers to the following questions. The expected format is a Word document, or similar, with a few images demonstrating your progress. The assignment will be submitted via Blackboard.
\begin{enumerate}
    \item What was the highest accuracy score you were able to achieve?
    \item Which features seemed to make the biggest impact in accuracy results?
    \item What troubles or difficulties did you encounter, and how did you resolve them?
    %\item What is the Rich Header?
\end{enumerate}


\end{document}